\section{Further questions}
\subsection{Joint optimization of arithmetic and graph sources}
A source with fewer coordinates need not yield the smallest height-one
polynomial: its monomial multiplicities, opposite-sign opportunities and
alias incidences also matter. Can arithmetic refactorization, row-sign
balancing and occurrence matching be chosen jointly with a useful
certificate of optimality for a fixed finite source? A rigorous lower
bound for one constrained compiler class would be meaningful; it would
still not be a lower bound for all Diophantine equations.

\subsection{Sharper pooled digit regions}
The present bounds exploit generous known radix margins. One may ask
whether an ellipsoidal form with certified integral weights permits further
pooling across unequal denominators without increasing degree or adding
more coordinates than it removes. Soundness must still identify the actual
integer remainder, including negative candidates, before factorial or
range decoding. A numerical search for small examples is not enough.

\subsection{Less expensive exact polynomial shears}
An integral coordinate shear can lower degree while increasing graph cost.
Can one share its polynomial-part calculation with the existing carrier
registers to retain the degree gain using fewer live binary operations?
The relevant objective is the full fixed-coefficient output graph, including
cancellations in the highest homogeneous component, not the apparent degree
of an intermediate expression.

\subsection{Formalizing the transformation boundary}
The new finite row transformations, monic witness restorations and network
energy identities are natural targets for machine-checked formalization.
A useful next milestone is a checked theorem connecting the emitted arrays
to the already identified source relation, followed by a transitive build
of the calibration boundary. Merely checking a supplied hash manifest or
including a named Lean declaration does not accomplish that milestone.
